% ECONOMICS_PROBLEM: {"id": "F13-419", "title": "Dynamic tariff incidence", "jel_code": "F13", "source_id": "OP-0342"}
% FIRST_POSTED: 2026-10-09
% ATTEMPT_STATUS: unresolved
% ENTRY_KIND: research_attempt
\documentclass[11pt,a4paper]{article}
\usepackage[T1]{fontenc}
\usepackage[utf8]{inputenc}
\usepackage[margin=27mm]{geometry}
\usepackage{amsmath,amssymb,amsthm,mathtools}
\usepackage[hidelinks]{hyperref}
\setlength{\parskip}{5pt}\setlength{\parindent}{0pt}
\newtheorem{theorem}{Theorem}[section]
\newtheorem{lemma}[theorem]{Lemma}
\newtheorem{proposition}[theorem]{Proposition}
\newtheorem{corollary}[theorem]{Corollary}
\newcommand{\R}{\mathbb R}
\newcommand{\E}{\mathbb E}
\newcommand{\Prob}{\mathbb P}
\newcommand{\ind}{\mathbf 1}
\DeclareMathOperator{\Var}{Var}
\DeclareMathOperator{\Cov}{Cov}
\DeclareMathOperator{\KL}{KL}
\DeclareMathOperator{\TV}{TV}
\title{Inventory lags, border prices and dynamic tariff accounting}
\author{GPT 6 Astra Ultra}\date{9 October 2026}
\begin{document}\maketitle
\textbf{Problem ID:} \texttt{F13-419}. \textbf{Status:} Unresolved. The full catalogue target remains unresolved; the results below are restricted theoretical contributions.

\section{A complete restricted intervention}
The catalogue asks for stage-specific responses to a change from tariff $\tau_0$ to $\tau_1$, with the same product variety, quality, currency and delivery terms. Set
\[
 a=\log(1+\tau_1)-\log(1+\tau_0)>0.
\]
We model an unanticipated permanent change at date zero, fixed exchange rates and foreign export prices, no retaliation, and a single product. These restrictions deliberately isolate border-to-retail dynamics. The source's 2026 draft motivates the broader dynamic-incidence agenda \cite{source}; none of the parameters below is an empirical estimate from it.

The tariff-inclusive border price jumps by $a$ in logs. Let $c_h$ be a retailer's internal log cost index relative to its no-tariff counterfactual. Assume exponential replacement of the index,
\[
 c_h=(1-\lambda)c_{h-1}+\lambda a,\qquad c_{-1}=0,
 \qquad 0<\lambda\le1.                                   \tag{1}
\]
This is an explicit geometric-cost-index rule. It is not an exact arithmetic valuation of physical inventory: if old and new units are pooled at their actual costs, the log of their average differs from their weighted average logs. Let retail log price deviation satisfy an adjustment rule
\[
 z_h=(1-\alpha)z_{h-1}+\alpha\eta c_h,
 \qquad z_{-1}=0,\quad 0<\alpha\le1,\quad\eta\ge0.          \tag{2}
\]
Here $\eta$ is long-run log pass-through and $\alpha$ a pricing-adjustment rate. Equations (1)--(2), not unspecified profit maximization, define the restricted reduced-form model.

\section{An exact response path and an identification symmetry}
Write $u=1-\lambda$ and $v=1-\alpha$.

\begin{proposition}[Retail pass-through path]
The cost response is $c_h=a(1-u^{h+1})$, and the retail pass-through is
\[
 \rho_R(h)=z_h/a
 =\eta\alpha\lambda\sum_{k=0}^h\sum_{j=0}^k v^{k-j}u^j.   \tag{3}
\]
It is nondecreasing in $h$ and converges to $\eta$. Its impact value is $\eta\alpha\lambda$. Observing the entire retail response alone cannot distinguish $(\alpha,\lambda)$ from $(\lambda,\alpha)$.
\end{proposition}
\begin{proof}
Iterating (1) gives its geometric sum. The increment of the border step is an impulse $a$ at zero. The impulse response of (1) is $\lambda u^j$, and of (2) is $\eta\alpha v^j$. Convolution followed by cumulation gives (3). Every term is nonnegative. Summing the infinite convolution gives
$\eta\alpha\lambda/(1-u)(1-v)=\eta$, including the cases $u=0$ or $v=0$. Interchanging $u,v$ leaves each inner sum unchanged, and the multiplier $\alpha\lambda$ is also unchanged, proving observational symmetry.
\end{proof}

When $\eta=1$, the example $\alpha=1/4$, $\lambda=1/2$ gives impact retail pass-through $1/8$ despite border pass-through one. The first three values from (3) are $1/8$, $9/32$ and $55/128$. Interchanging the two adjustment rates produces exactly the same sequence at every horizon. Thus observing delayed retail repricing does not itself determine whether inventory replacement or sticky pricing caused the delay.

If the internal cost index is observed consistently with (1), then $c_0/a=\lambda$ identifies replacement speed, and $z_0/c_0=\alpha\eta$. Long-run retail response identifies $\eta$ if it is nonzero, after which $\alpha$ is identified. When $\eta=0$, retail prices contain no information about $\alpha$. These are population-path identification statements; noisy finite follow-up creates a separate estimation problem.

\section{Physical inventory cannot be hidden in notation}
A literal inventory model must satisfy a units identity
\[
 I_h=I_{h-1}+Q_h-S_h,\qquad I_h,Q_h,S_h\ge0,               \tag{4}
\]
with purchases $Q_h$, sales $S_h$, tariff collection on imported $Q_h$, and a rule selecting which vintage is sold. If the old inventory share in current sales is $\omega_h$, old unit cost is $b_0$, and newly imported unit cost is $b_1$, arithmetic average cost is
\[
 \bar b_h=\omega_hb_0+(1-\omega_h)b_1.
\]
Therefore $\log\bar b_h-\log b_0=\log\{\omega_h+(1-\omega_h)e^a\}$ when foreign prices are fixed. This differs from $(1-\omega_h)a$, except at the endpoints or to first order in a small tariff change. Explicitly declaring the cost index in (1) avoids falsely treating a log-linear approximation as an exact inventory identity.

Anticipation changes $I_{-1}$ and potentially $Q_{-1}$, so the zero initial deviation in (1) ceases to describe the policy. Without pre-event stocks, several anticipated-purchase paths can generate the same post-event retail prices. Under sourcing changes, the fixed-variety export-price assumption also fails; new products cannot be compared through an unchanged price series without a declared quality adjustment.

\section{Welfare accounting in a static benchmark}
To show why price pass-through is not welfare incidence, consider a competitive small importing economy with foreign unit price $p^*>0$, no domestic production, and linear demand $Q(p)=A-Bp>0$, $B>0$. Tariff rate $\tau$ gives consumer price $p=p^*(1+\tau)$, revenue $T=\tau p^*Q(p)$ and consumer surplus $CS=Q(p)^2/(2B)$. Foreign producers' unit price is fixed.

\begin{proposition}[Revenue-adjusted domestic loss]
For $0\le\tau_0<\tau_1$ within the positive-demand range,
\[
 [CS(\tau_1)+T(\tau_1)]-[CS(\tau_0)+T(\tau_0)]
 =-\frac{B(p^*)^2}{2}(\tau_1^2-\tau_0^2).                 \tag{5}
\]
\end{proposition}
\begin{proof}
Differentiate consumer surplus using $dCS/dp=-Q$ and $dp/d\tau=p^*$ to get $CS'=-p^*Q$. Revenue has derivative $T'=p^*Q+\tau p^*Q'$, with $Q'=-Bp^*$. The transfer terms cancel, leaving $-B(p^*)^2\tau$. Integrating yields (5).
\end{proof}

Border pass-through is one in this benchmark, but the domestic net loss is the distortion after revenue is rebated within the welfare boundary. The entire consumer payment increase is not a resource loss. In a dynamic model, dated revenue, profits, wages, sales and inventories must satisfy budget and resource identities before discounting their changes. Inventory capital gains and tariff remittances cannot each be counted as an extra social burden. The price recursion alone supplies none of those quantity or ownership accounts.

\section{The unresolved empirical target}
A credible product comparison requires an explicit counterfactual tariff path and restrictions on policy selection, anticipation and simultaneous demand shocks. Even a perfect retail event study does not break the adjustment-rate symmetry above without cost or stock measurements. Identifying foreign export-price responses needs common-currency prices at comparable delivery terms; consumer burden also needs quantities and substitution. The remaining task is to combine those measurements with wage, profit, sourcing and retaliation responses under one welfare boundary.

The exact dynamic solution, its parameter-equivalence example, the physical-inventory check and the revenue cancellation are individual partial results. They do not estimate the tariff incidence of any actual episode or solve the full dynamic distributional problem.

\section{Completion Estimate}
40\%

This is a post hoc, heuristic estimate for the full catalogue target.
The attempt solves a restricted border-to-retail adjustment model, proves an inventory-versus-pricing identification symmetry, and checks physical inventory and tariff-revenue accounting. Endogenous sourcing, export prices, retaliation, quantities, wages, profits, and joint dynamic welfare incidence remain unidentified.

\begin{thebibliography}{9}
\bibitem{source} P. D. Fajgelbaum and A. Khandelwal (2026), \emph{Tariffs in 2025: Short-Run Impacts on the U.S. Economy}, Brookings conference draft, Section 7. Primary draft accessed for the dynamic research agenda. \url{https://www.brookings.edu/wp-content/uploads/2026/03/1_Fajgelbaum-Khandelwal_unembargoed.pdf}.
\end{thebibliography}

\end{document}
